\documentclass[11pt]{article}
\usepackage[margin=1in]{geometry}
\usepackage{amsmath,amssymb,amsthm,mathtools,mathrsfs}
\usepackage{microtype}
\usepackage{enumitem}
\usepackage{booktabs}
\usepackage[hidelinks]{hyperref}
\setlist{nosep}
\newtheorem{theorem}{Theorem}[section]
\newtheorem{proposition}[theorem]{Proposition}
\newtheorem{corollary}[theorem]{Corollary}
\newtheorem{lemma}[theorem]{Lemma}
\theoremstyle{definition}
\newtheorem{definition}[theorem]{Definition}
\newtheorem{warning}[theorem]{Warning}
\theoremstyle{remark}
\newtheorem{remark}[theorem]{Remark}
\newcommand{\A}{\mathcal A}
\newcommand{\T}{\mathcal T}
\newcommand{\F}{\mathscr F}
\newcommand{\PP}{\mathbf P}
\newcommand{\QQ}{\mathbf Q}
\title{Proposed Manuscript Insertions VII\\
Source-Complete Binary Contacts and the Pseudo-Trace Atlas}
\author{Research draft for the marked-width MLC program}
\date{July 25, 2026}

\begin{document}
\maketitle

\begin{abstract}
For the actual binary $f^q$ pullback system at the first direct-decoration
scale, we recover the missing one-step boundary contacts from the source
$2^n$-component attachment theorem, compute all deeper contacts by a finite
normalized recursion, retain endpoint flags and cyclic sides, and construct a
one-to-many trace-pullback atlas on the source-native pseudo-puzzle boundary.
The atlas is exactly natural under normalization, safe-gap carrier enlargement,
and certified finite overlay refinement.  The results close the binary
contact/trace construction for a fixed source block and finite marked overlay;
they do not prove a universal response/Value quotient, MBTT, rigidity, or MLC.
\end{abstract}

\section{Exact binary source package}
Let
\[
 P=Y^{(n-1)q+1},\qquad g=f^q,
\]
and let $Q_L,Q_R\subset P$ be the two actual univalent $f^q$ pullbacks.
Write
\[
 g_a=g|_{Q_a}:Q_a\to P,\qquad \phi_a=g_a^{-1},\qquad a\in\A=\{L,R\}.
\]
For $w=a_1\cdots a_k$, put
\[
 \phi_w=\phi_{a_1}\circ\cdots\circ\phi_{a_k},\quad
 D_w=\phi_w(P),\quad h_w=g^k|_{D_w}:D_w\to P.
\]
The already established closed-piece separation gives $Q_L\cap Q_R=\varnothing$.

The source proves that the domain of $g^n$ has $2^n$ components.  For
$s\in\A^n$, let
\[
 v_s=\phi_s(\alpha).
\]
These points are the vertices of $P$, and the component $D_s$ is attached to
the two boundary rays of $P$ landing at $v_s$.

\section{One-step contact recovery}
Define
\[
 C_a=\partial P\cap\partial Q_a,
 \qquad
 \mathcal V_a=\{v_{at}:t\in\A^{n-1}\}.
\]

\begin{theorem}[Source recovery of the binary contacts]
For $a\in\{L,R\}$:
\begin{enumerate}
\item the landing vertices of $C_a$ are exactly $\mathcal V_a$;
\item $C_a$ consists of the maximal common subsegments of the two boundary
rays of $P$ landing at those vertices, together with their endpoint flags;
\item no nondegenerate equipotential arc belongs to $C_a$;
\item the complete contact is computed by exact matching of the finite typed
boundary words of $P$ and $Q_a$.
\end{enumerate}
After finite subdivision, $g_a$ maps every open contact atom homeomorphically
onto one root-boundary atom.
\end{theorem}

\begin{proof}
For $s=at$, the source component $D_s$ is contained in $Q_a$ and shares the
pair of boundary rays of $P$ landing at $v_s$.  A point of such a common ray
cannot be interior to $Q_a\subset P$, hence the common ray is also on
$\partial Q_a$.  Thus every $v_{at}$ occurs.

Conversely, let an open contact ray of $P$ and $Q_a$ land at a vertex $v$ of
$P$.  The source attachment theorem supplies the unique depth-$n$ component
$D_s$ attached at $v$.  Its first-level cylinder contains $v$.  Since the two
closed first-level cylinders are disjoint and $v\in Q_a$, the first symbol of
$s$ must be $a$.  Hence $v=v_{at}$.

The parent and child equipotential levels differ by $q$, so their open arcs are
disjoint.  Distinct external rays are disjoint outside the filled Julia set.
Thus exact ray-angle matching and intersection of the potential intervals
recovers all open contacts.  Subdivision at their endpoints and images makes
the branch map simplicial. 
\end{proof}

\begin{definition}[Endpoint flag]
A boundary flag is a vertex, an incident open boundary atom, the side of the
puzzle piece, and the induced cyclic orientation.  Bare landing points are not
used as endpoint records.
\end{definition}

Each $g_a$ induces a partial flag map
\[
 \sigma_a:\F(C_a)\longrightarrow\F(\partial P)
\]
that preserves cyclic order.

\section{Complete all-depth contacts}
For a word $w$, set
\[
 C_w=\partial P\cap\partial D_w,
 \qquad
 \widehat C_w=h_w(C_w)\subset\partial P,
 \qquad
 \widehat C_{\varnothing}=\partial P.
\]

\begin{theorem}[Exact contact recursion]
For every word $w$ and symbol $a$,
\[
 \boxed{\widehat C_{wa}=g_a(\widehat C_w\cap C_a).}
\]
\end{theorem}

\begin{proof}
If $z\in C_{wa}$, nesting $D_{wa}\subset D_w\subset P$ implies
$z\in\partial D_w$: otherwise $z$ would be interior to $D_w\subset P$ while
lying on $\partial P$.  Hence $x=h_w(z)\in\widehat C_w$.  Since
$D_{wa}=\phi_w(Q_a)$, we also have $x\in\partial Q_a$, and
$h_{wa}(z)=g_a(x)$.

Conversely, if $x\in\widehat C_w\cap C_a$, then
$z=h_w^{-1}(x)$ lies on $\partial P$ by the definition of $\widehat C_w$, and
on $\partial D_{wa}$ because $x\in\partial Q_a$.  Thus $z\in C_{wa}$.
\end{proof}

\begin{theorem}[Complete binary cylinder contacts]
For finite words $u,v$:
\begin{enumerate}
\item if $u,v$ are incomparable, then $D_u\cap D_v=\varnothing$;
\item if $v=ut$, then
\[
 \boxed{\partial D_u\cap\partial D_v=\phi_u(C_t).}
\]
\end{enumerate}
Consequently the two one-step contact graphs determine every all-depth contact.
\end{theorem}

\begin{proof}
At the first differing symbol, incomparable pieces lie in the inverse images of
the disjoint closed pieces $Q_L,Q_R$.  If $v=ut$, then
$D_v=\phi_u(D_t)$ and $D_u=\phi_u(P)$; apply the closed-piece homeomorphism
$\phi_u$.
\end{proof}

After subdivision, the normalized contact dynamics is the finite automaton
\[
 T_a(S)=g_a(S\cap C_a)
\]
on subcomplexes of one finite root boundary graph.  This is a per-block finite
state statement, not a universal bound in $n$ or $q$.

\section{Source endpoint itineraries}
For a prefix $w\in\A^k$, $0\le k\le n$, define
\[
 \mathcal V(w)=\{v_{wt}:t\in\A^{n-k}\}.
\]

\begin{theorem}[Prefix endpoint itinerary]
The landing vertices of the open ray contacts in
$\partial P\cap\partial D_w$ are exactly $\mathcal V(w)$.  Moreover
\[
 \boxed{h_w(v_{wt})=v_t.}
\]
The incident boundary-ray flags and their cyclic order are transported by the
same shift.
\end{theorem}

\begin{proof}
Every depth-$n$ component $D_{wt}\subset D_w$ is source-attached to the
boundary rays at $v_{wt}$; nesting forces those common rays to remain on the
boundary of each intermediate piece.  Conversely, a contact vertex of $D_w$
has a unique source depth-$n$ attached component.  Its word must begin with
$w$, since otherwise it lies in an incomparable disjoint cylinder.  Finally,
$g^k$ cancels the first $k$ inverse branches.
\end{proof}

\begin{theorem}[Post-transient synchronization]
Let $\mathscr R_\alpha\subset\partial P$ be the complete pair of root
boundary-ray atoms landing at $\alpha$, with endpoint flags.  For every
$s\in\A^n$,
\[
 \widehat C_s=\mathscr R_\alpha.
\]
The left branch fixes this contact, while the right branch kills it:
\[
 T_L(\mathscr R_\alpha)=\mathscr R_\alpha,
 \qquad
 T_R(\mathscr R_\alpha)=\varnothing.
\]
Hence for every continuation $t$, the contact persists exactly when
$t=L^m$; after the first $R$ it is empty forever.  The same holds for endpoint
flags.
\end{theorem}

\begin{proof}
The prefix theorem at $k=n$ gives the unique normalized landing vertex
$\alpha$ for every depth-$n$ word, and $h_s$ maps the attached deeper ray pair
onto the complete root ray pair.  The left inverse branch is attached at the
fixed point $\alpha$ and $f^q$ fixes its ray pair.  The right first-level piece
is a disjoint closed piece attached at $\alpha'$, so exact boundary-word
matching gives no ray atom or endpoint flag in common with
$\mathscr R_\alpha$.  Apply the contact recursion; the empty state is
absorbing.
\end{proof}

\begin{corollary}
The $2^n$ source fan is transient contact data.  Beyond source depth $n$, one
address has only two possible normalized contact states: the fixed
$\alpha$-port or no root contact.  This does not bound the number of transient
proof-used ports across all words.
\end{corollary}

\section{One-to-many pseudo-trace atlas}
The source pseudo-functor associates a pseudo-piece $\mathbf Y$ to each puzzle
piece.  Its ideal boundary has finitely many O-boundary and J-boundary arcs; at
a vertex $v$ the source port is $\partial_v\mathbf Y$.  A degree-one puzzle
branch induces a degree-one covering of pseudo-pieces and extends continuously
to the ideal boundary.

Choose a finite root overlay $\mathscr O$ containing all contact endpoints,
source vertex-port endpoints, and the current finite protected, obstacle and
electrode endpoints.  Pull it back to $\mathbf Q_L,\mathbf Q_R$ and refine the
root by all images.  Let $\mathscr A_P$ and $\mathscr A_{Q_a}$ be the resulting
open atoms.  Each child atom $K$ maps homeomorphically to one parent atom
$p_a(K)$.

For $J\in\mathscr A_P$, put
\[
 \operatorname{Lift}(J)=\{(a,K):p_a(K)=J\}.
\]
Because each branch has degree one,
\[
 |\operatorname{Lift}(J)|\le2.
\]

For a pseudo-piece $\mathbf Y$ and a finite atom union $B$, let
\[
 \T_{\mathbf Y}(B)
 =H^1_{\mathrm{fin}}(\mathbf Y)/\{u:\operatorname{Tr}_B u=0\}
\]
with its minimum-energy quotient norm.  Define the compatible tuple space over
$J$ to be the image of the joint branch pullback.

\begin{theorem}[Binary one-to-many trace atlas]
For every root atom $J$, actual trace pullback defines a continuous map
\[
 \mathcal R_J:\T_{\PP}(J)\longrightarrow
 \T_J^{\mathrm{comp}},
\]
\[
 \mathcal R_Ju=
 \left((u|_{\mathbf g_a(K)})\circ\mathbf g_a|_K\right)_{(a,K)\in\operatorname{Lift}(J)}.
\]
It retains exact branch pullback, endpoint flags, side and cyclic order; it is
conformally natural for Dirichlet energy and assumes no omitted Euclidean
boundary welding.
\end{theorem}

\begin{proof}
The source-induced pseudo-map is conformal and extends to the ideal boundary.
Every refined child atom maps homeomorphically to its parent atom.  Therefore
trace pullback is defined on the intrinsic minimum-energy trace quotient.  The
compatible tuple space is the image of the joint map, so vertex compatibility
is automatic.  Conformal invariance gives energy naturality.
\end{proof}

For a word $w$, direct pullback by the induced $\mathbf h_w$ equals the
composition of the one-step maps, with canonical reassociation of the
compatible tuple spaces.  The analogous flag and cyclic-order identities also
hold.

\section{Normalized-overlay naturality}
At a word $w$, put
\[
 \mathscr O_w=\mathbf h_w^{-1}(\mathscr O).
\]
On $\mathbf D_{wa}$,
\[
 \boxed{\mathbf h_{wa}=\mathbf g_a\circ\mathbf h_w.}
\]

\begin{theorem}[Normalized-overlay naturality]
For every $w,a$, this identity gives commuting diagrams for:
\begin{enumerate}
\item contact subgraphs;
\item endpoint flags and cyclic order;
\item one-to-many trace restrictions;
\item compatible trace spaces;
\item all finite response forms obtained by minimizing internal traces.
\end{enumerate}
Two prefixes with the same normalized contact state and live marked overlay
therefore have the same future contact/trace behavior.
\end{theorem}

\begin{proof}
The polynomial identity is exact on the selected pieces.  The pseudo-functor
preserves it.  The contact square is the contact-recursion theorem; the flag
and trace squares are the boundary extension and functoriality of restriction.
Dirichlet minimization is taken over coordinate-isomorphic fibers.
\end{proof}

\section{Final-root-enlargement naturality}
\begin{theorem}[Safe-gap atlas invariance]
Let
\[
 K_0\subset S\subset S^+\subset P\setminus(Q_L\cup Q_R)
\]
be compact connected roots, and put $S_w=\phi_w(S)$,
$S_w^+=\phi_w(S^+)$.  Then every puzzle piece, contact set, endpoint itinerary,
flag, and trace-atlas map is unchanged.  The inclusions $S_w\subset S_w^+$
commute with all shift and parent maps; only the internal carrier/cap response
can change.
\end{theorem}

\begin{proof}
The atlas depends only on the fixed puzzle and pseudo-puzzle boundaries and on
the finite overlay.  The safe-gap enlargement does not meet them.  The
inclusions and commuting squares follow from injectivity of $\phi_w$ and
$g\circ\phi_a=\mathrm{id}$.
\end{proof}

For a certified finite overlay refinement
$\mathscr O\preceq\mathscr O^+$, assume every new child endpoint is an exact
pullback or typed source endpoint, every branch is atomwise homeomorphic, and
all protected sides are retained.  Let $\pi_P,\pi_a$ forget the subdivision.

\begin{theorem}[Refinement naturality]
The atom, flag and trace diagrams commute:
\[
 \pi_Pp_a^+=p_a\pi_a,
 \qquad
 \pi_P\sigma_a^+=\sigma_a\pi_a,
\]
\[
 \boxed{\operatorname{Sub}_{Q_a}\mathcal R_a
 =\mathcal R_a^+\operatorname{Sub}_P.}
\]
Introducing and then minimizing the refined internal traces reconstructs the
coarse response exactly.
\end{theorem}

\begin{remark}[Promotion boundary]
This proves naturality of a certified collision-free final overlay.  It does
not prove that an arbitrary final-root enlargement is collision-free or
proof-complete.  No-new-demand stability still requires demanded-atom
bijection or future observability on every coarse fiber.
\end{remark}

\section{Recursive response grammar and finite-state observables}
Let $\mathcal B$ be the exact one-step Dirichlet form on the root terminals and
the two child trace interfaces, and let $Q_0$ be a terminal cap response.  With
$\star$ denoting exact compatible-trace gluing, define
\[
 Q_{r+1}=\mathcal B\star\bigl(Q_r^{(L)}\oplus Q_r^{(R)}\bigr).
\]

\begin{theorem}[Binary response grammar]
$Q_r$ is exactly the response of the normalized depth-$r$ binary assembly.
It has a shared recursive circuit of size $O(r)$, although its fully expanded
leaf interface may have $2^r$ ports.  The recursion is natural under the
normalization and enlargement maps above.
\end{theorem}

\begin{proof}
The depth-$r+1$ domain is one root cell with one normalized depth-$r$ assembly
on each child.  Apply Sobolev domain decomposition and induction.  Sharing the
identical child response symbol gives one new circuit node per depth.
\end{proof}

Let $M_L,M_R$ be $d\times d$ coefficient transitions and $R_L,R_R$ be
$w\times w$ branch-feature transports.  For $s=a_1\cdots a_r$, write
$M_s=M_{a_1}\cdots M_{a_r}$ and $R_s=R_{a_1}\cdots R_{a_r}$.

\begin{theorem}[Matrix-product future quotient]
For coefficient vectors $u,v$ and feature vector $x$,
\[
 \sum_{s\in\{L,R\}^r}(u^TM_sv)R_sx
 = (u^T\otimes I_w)
   (M_L\otimes R_L+M_R\otimes R_R)^r
   (v\otimes x).
\]
Hence a universal $d$-state coefficient transducer and $w$-dimensional branch
feature module give an exact future quotient of dimension at most $dw$,
independent of $2^r$.
\end{theorem}

\begin{proof}
Expand the $r$-th power.  Its terms are exactly $M_s\otimes R_s$, indexed by
binary words, and contract the first tensor factor with $u,v$.
\end{proof}

\subsection*{Fixed-block KKL Value moment}
Let $T$ be one source-licensed finite terminal Hubbard tree, $E=E(T)$, and
$\mathcal S=E\times E$.  Let $v,w$ be the Perron weights and $\lambda$ the core
eigenvalue used by the KKL Value construction.  For a level-$m$ arc $\gamma$,
define
\[
 z_m(\gamma)_{e,e'}=
 \sum_{\substack{\alpha:\,f_*^m\alpha=e\\
                 (\iota^m)_*\alpha=e'}}
 \langle\alpha,\gamma\rangle,
 \qquad
 a_{e,e'}=v(e)\sqrt{w(e')}.
\]
For a WAD $X$, put
\[
 \mathsf M_m(X)=\sum_\gamma X(\gamma)
 z_m(\gamma)z_m(\gamma)^T.
\]

\begin{theorem}[Exact fixed-block Value state]
\[
 \boxed{
 \operatorname{Val}(X)=\lambda^{-2m}
 \langle aa^T,\mathsf M_m(X)\rangle.}
\]
Thus one fixed terminal tree has an exact positive homogeneous Value state of
dimension at most $|E|^2(|E|^2+1)/2$.  If a univalent binary branch induces
an incidence matrix $P_a$ on $\mathbb R^{\mathcal S}$, then
\[
 \mathsf M(b_a^*X)=P_a\mathsf M(X)P_a^T,
\]
and the joint binary transport is
\[
 \mathsf M\longmapsto
 P_L\mathsf MP_L^T+P_R\mathsf MP_R^T.
\]
\end{theorem}

\begin{proof}
The KKL coefficient of a pulled bivertical arc $\alpha$ is
$\lambda^{-m}v(f_*^m\alpha)\sqrt{w((\iota^m)_*\alpha)}$.
Grouping by the pair $(f_*^m\alpha,(\iota^m)_*\alpha)$ gives
$\ell_m(\gamma)=\lambda^{-m}a^Tz_m(\gamma)$.  Substitute this into
$\operatorname{Val}(X)=\sum_\gamma X(\gamma)\ell_m(\gamma)^2$.
Univalent branch transport preserves intersection numbers, yielding the
congruence formula.
\end{proof}

\begin{warning}
The matrix-product theorem and the fixed-block moment theorem do not provide a
universal state by themselves.  The terminal edge dimension may grow, and a
virtual block still needs an actual typed factor to the terminal bivertical
incidence.  Arbitrary independent leaf coefficients still require exponentially
many degrees of freedom.
\end{warning}

\section{Status and next gate}
The following fixed-block binary modules are proved by the preceding arguments:
source contact recovery, complete contact recursion, complete cylinder
contacts, endpoint-germ itineraries, the one-to-many pseudo-trace atlas, word
functoriality, normalized-overlay naturality, safe-gap atlas invariance, and
certified-refinement naturality.

The broad direct-decoration and MLC program still requires a complete source
Wave horizon, protected final inventory, a universal future-observable
response/WAD/Value quotient, MBTT, coercivity and rigidity.  The next major
analytic theorem should construct a uniformly controlled binary trace/Value
future quotient and a strict complete-block transition.  No MLC claim is made.

\end{document}
